\documentclass[11pt]{article}

\usepackage[margin=1in]{geometry}
\usepackage[T1]{fontenc}
\usepackage[utf8]{inputenc}
\usepackage{lmodern}
\usepackage{microtype}
\usepackage{graphicx}
\usepackage{booktabs}
\usepackage{amsmath,amssymb,amsthm}
\usepackage[numbers,sort&compress]{natbib}
\usepackage[colorlinks=true,allcolors=blue]{hyperref}

\title{A Lightweight Baseline for Structured Prompt Auditing}
\author{Noramiya et al.}
\date{\today}

\begin{document}

\maketitle

\begin{abstract}
This sample demonstrates a compact, publication-style manuscript generated from the \texttt{codex-latex} skill.
The paper structure enforces stable section ordering, deterministic labels, and consistent bibliography keys.
A simple auditing baseline is used to show how methods and results can be documented with minimal formatting overhead.
The template is suitable for early-stage research drafting and can be upgraded to venue-specific classes later.
\end{abstract}

\section{Introduction}
Research teams that use AI assistants often need fast paper drafts with reliable formatting.
We build a baseline workflow that focuses on predictable structure and easy maintenance, following long-standing LaTeX practices \citep{lamport1994}.
Our goal is not to optimize writing quality directly, but to make the authoring process reproducible.

\section{Method}
The workflow has three steps: (1) scaffold a project with a selected template, (2) edit sections with stable labels and citation keys, and (3) run compile and reference checks before delivery.
Equation~\eqref{eq:score} defines a simple audit score used in this sample.

\begin{equation}
\label{eq:score}
S = \frac{n_{\text{resolved}}}{n_{\text{total}}}
\end{equation}

\section{Results}
Table~\ref{tab:results} shows an example summary.
The baseline improves reference consistency while keeping document edits small and reviewable.

\begin{table}[t]
  \centering
  \caption{Illustrative metrics for the sample workflow. Higher is better.}
  \label{tab:results}
  \begin{tabular}{lcc}
    \toprule
    Setup & Consistency Score & Edit Churn \\
    \midrule
    Ad-hoc drafting & 0.71 & High \\
    Structured skill workflow & 0.94 & Low \\
    \bottomrule
  \end{tabular}
\end{table}

\begin{figure}[t]
  \centering
  \fbox{\rule{0pt}{1.2in}\rule{0.9\linewidth}{0pt}}
  \caption{Placeholder figure block used to keep the sample self-contained.}
  \label{fig:placeholder}
\end{figure}

Figure~\ref{fig:placeholder} demonstrates a shareable figure slot without external assets.

\section{Discussion}
This sample emphasizes maintainability over novelty.
In practical projects, teams can replace placeholders with real experiments while preserving labels, citations, and section semantics.

\section{Conclusion}
The generated manuscript provides a clean baseline for research writing workflows.
It is immediately shareable and can be adapted to journal or conference templates with minimal migration effort.

\bibliographystyle{plainnat}
\bibliography{references}

\end{document}
